\section{Method}
\label{sec:method}

\subsection{Study Design}
This study examines whether post-training quantization to 8-bit integer precision
(INT8) causes a larger relative detection-accuracy loss for small objects than
for larger objects in aerial imagery. Precision is the experimental factor, with
a TensorRT FP32 engine as the baseline, an INT8 engine as the primary treatment,
and an FP16 engine as a control. Object size is an observed moderator. The
primary comparison is paired: all precision conditions are evaluated on the same
images.

For each object-size group $g$, relative AP loss is defined as
\begin{equation}
\mathrm{RelDrop}_{g} =
\frac{\mathrm{AP}^{\mathrm{FP32}}_{g}
- \mathrm{AP}^{\mathrm{INT8}}_{g}}
{\mathrm{AP}^{\mathrm{FP32}}_{g}}.
\label{eq:relative-ap-drop}
\end{equation}
The primary contrast is
\begin{equation}
D = \mathrm{RelDrop}_{\mathrm{small}}
- \mathrm{RelDrop}_{\mathrm{non\text{-}small}}.
\label{eq:size-interaction}
\end{equation}
The alternative hypothesis is that $D>0$, indicating a greater relative
quantization loss for small objects. The null hypothesis is $D\leq 0$.

\subsection{Dataset and Annotation Processing}
The experiment uses VisDrone2019-DET, with the training split used for model
fine-tuning and INT8 calibration, the validation split used for checkpoint
selection and development, and test-dev reserved for final evaluation. The
preregistered split sizes are 6,471 training images, 548 validation images, and
1,610 test-dev images. Dataset archives are verified by SHA-256 before
processing.

Annotations are converted to YOLO format for training and to COCO JSON for
evaluation. VisDrone annotations with score zero, including ignored regions and
the ``others'' category, are excluded. Bounding boxes are clipped to the image
boundaries, and boxes with a width or height below one pixel after clipping are
discarded. Consequently, detections on excluded ignored regions are treated as
false positives, and the resulting absolute AP values are not directly
comparable with scores from the official VisDrone evaluation toolkit.

Object size is calculated from the clipped ground-truth box area in original
image pixels. Small objects have area below $32^2=1{,}024$ pixels; non-small
objects have area at least $1{,}024$ pixels. Secondary size groups are medium
($32^2 \leq \mathrm{area} < 96^2$) and large
($\mathrm{area} \geq 96^2=9{,}216$ pixels). Size-specific evaluation follows
the COCO convention: ground truths outside a size range, detections matched to
those ground truths, and unmatched detections outside the range are ignored for
that range.

\subsection{Model Training and Engine Construction}
The detector is YOLOv8n initialized from COCO-pretrained weights. One model is
fine-tuned on the training split at an image size of $640 \times 640$ for 100
epochs, with batch size 16, random seed 0, and deterministic training enabled.
Other training options use the Ultralytics defaults; no hyperparameter search is
performed. The checkpoint with the best validation AP is frozen and used to
construct all precision variants.

The frozen checkpoint is exported once to ONNX with a static input shape of
$1 \times 3 \times 640 \times 640$. TensorRT engines are then built in FP32,
FP16, and INT8. Three independent INT8 calibration sets are sampled, one per
seed (0, 1, and 2); each contains 512 images sampled uniformly from the training
split. Calibration and evaluation use the same $640 \times 640$ letterbox
preprocessing. The primary INT8 comparison uses calibration seed 0; seeds 1 and
2 are reserved for robustness analysis. The TensorRT 11 toolchain uses explicit
quantize/dequantize quantization through NVIDIA ModelOpt. If the preregistered
TensorRT 10 fallback is required, its implicit calibration path is recorded as
a deviation from the preferred toolchain.

\subsection{Evaluation Measures}
The frozen PyTorch checkpoint and each available TensorRT engine are evaluated
using identical inference settings: confidence threshold 0.001, non-maximum
suppression (NMS) intersection-over-union threshold 0.7, a maximum of 500
detections per image, batch size 1, and non-rectangular inference at image size
640. The primary accuracy measure is COCO-style AP averaged over IoU thresholds
0.50 through 0.95 (AP@[.5:.95]); AP at IoU 0.50 (AP@0.5) is secondary.
Evaluation is reported overall and by the predefined object-size groups.

Predictions are matched to ground truth using pycocotools, with at most 500
detections per image. The evaluation implementation also stores the per-image
match records used in the bootstrap analysis. Its AP recomputation is checked
against pycocotools and must agree within an absolute difference of $10^{-9}$.
As implementation checks, the PyTorch FP32 checkpoint is compared with the
TensorRT FP32 engine, and the custom evaluation results are checked against
Ultralytics validation metrics.

All development and debugging evaluations use validation data. The test-dev
split is reserved for the final evaluation and is evaluated only after the
pipeline and analysis procedure are complete. The confirmatory comparison is
TensorRT FP32 versus INT8 calibration seed 0 on test-dev.

\subsection{Statistical Analysis}
The unit of resampling is the image. A paired nonparametric bootstrap draws
2,000 resamples of the 1,610 test-dev images with replacement, using random
seed 0. The same resampled image indices are used for every precision condition
and contrast. For each resample, AP is recomputed by object-size group from the
stored match records; detections are not re-matched.

The confirmatory statistic is $D$ in \eqref{eq:size-interaction} for the
FP32-versus-INT8-seed-0 comparison using AP@[.5:.95]. A two-sided 95\%
percentile bootstrap confidence interval is calculated for $D$. If its lower
bound is above zero, the null hypothesis is rejected in favor of a greater
relative loss for small objects. If the interval contains zero, the result is
inconclusive; if its upper bound is below zero, the result is contrary to the
alternative hypothesis.

All other comparisons are exploratory and are reported with confidence
intervals without additional hypothesis tests. They include absolute AP drops,
AP@0.5, small/medium/large AP, INT8 calibration seeds 1 and 2, and the FP16
control. The INT8 result is considered directionally robust if $D$ has the same
sign for all three calibration seeds. FP16 is used to interpret whether a
size-dependent effect is specific to INT8: a similar pattern under FP16 would
prevent attributing the effect to INT8 alone.

\subsection{Inference Latency}
Latency is a secondary descriptive measure, not part of the confirmatory
hypothesis test. The preregistered benchmark measures one engine forward call
through Ultralytics \texttt{AutoBackend} using a preprocessed
$1 \times 3 \times 640 \times 640$ tensor already resident in GPU memory.
Preprocessing, host-to-device transfer, decoding, and NMS are excluded. GPU
execution is synchronized immediately before and after timed calls, which are
measured with \texttt{time.perf\_counter}. FP32 and INT8 seed 0 are measured in
three alternating-order rounds; each round has 200 warm-up calls followed by
1,000 timed calls. The report is to include the median, interquartile range,
per-round medians, and the ratio of the FP32 and INT8 medians as the speed-up.
No statistical tests are planned for latency. This benchmark is reported only
if it is run; otherwise latency is omitted and the omission is documented.
